% =====================================================================
%  §4  A∞-algebras and A∞-categories  （原文 p22–p28）
% =====================================================================
\section{$A_\infty$-代数与 $A_\infty$-范畴}

\subsection{一般定义的两个问题}

本节的目的是描述一个框架，关于 $A_\infty$-范畴的结果将在其中表述。在回顾 Fukaya 范畴的定义之前，
我们想先作几点评论。这些评论是非正式的。精确的定义将在本节后面给出。

Fukaya 范畴的定义有两个主要问题。第一，态射只能对横截（transversal）的拉格朗日子流形定义
（特别地，单位态射永远没有定义）。第二，由于存在以给定拉格朗日子流形为边界的伪全纯圆盘，
必须在复合 $m_n$（$n\ge1$）之外再添上一个复合 $m_0$。其结果是，态射空间不再是复形：
$m_1^2\ne0$。另一方面，凝聚层的导出范畴来自一个没有 $m_0$、且满足 $m_1^2=0$ 的
$A_\infty$-范畴。因此必须说明，所讨论的两个 $A_\infty$-范畴在何种意义下等价。上述问题可以
通过适当推广 $A_\infty$-范畴的概念来解决。我们在下面给出这样一种推广（称之为
$A_\infty$-预范畴，$A_\infty$-pre-category）。我们要指出，还存在一个更好的推广。它需要
大量准备工作，将另文给出（见 \cite{KoS}）。另一方面，在本文所考虑的阿贝尔簇情形的 Fukaya
范畴版本中，$m_0$ 带来的问题并不出现。因此，就本文的目的而言，使用 $A_\infty$-预范畴
（或带横截结构的 $A_\infty$-范畴，参看 \cite{P1}）已经足够。这给出了横截性问题的一个部分
解决，同时给出了单位态射问题的一个解决。

我们将利用 $A_\infty$-预范畴，在 \S8 中表述并证明环面纤维化的同调镜像猜想的一个变体。
它可以应用于阿贝尔簇的情形。特别地，由此可以得到阿贝尔簇上 Massey 积的某些公式，其表达
形式是类似于 \cite{P1} 中所考虑的部分 theta 和（partial theta-sum）。

\begin{yizhu}[「两个问题」指什么]
原文标题说「两个问题」，正文随即点明：其一是态射只对横截的拉格朗日子流形对可定义，因而
连单位态射都不存在；其二是必须引入 $m_0$，因而 $m_1^2\ne0$、态射空间不是复形。本节后面
的两个推广正分别针对这两点：非幺 $A_\infty$-预范畴（定义 7）让 $\Hom$ 只在横截对上取定义，
$A_\infty$-预范畴的扩张性质（定义 10）则用「同伦意义下的单位」替代单位态射。
\end{yizhu}

\subsection{非幺 $A_\infty$-代数与 $A_\infty$-范畴}

设 $A=\bigoplus_{i\in\Z}A^i$ 是域 $k$ 上的 $\Z$-分次模。\footnote{今后可以把 $k$ 换成一个
$\Z$-分次交换结合代数，并假定所有 $k$-模都是投射模。}照例，用 $A[n]$ 记满足
$(A[n])^i=A^{i+n}$（对一切 $i$）的分次 $k$-模。

\begin{dfn}
$A$ 上的非幺 $A_\infty$-代数（non-unital $A_\infty$-algebra）结构，是指在余自由张量余代数
（cofree tensor coalgebra）$T_+(A[1])=\bigoplus_{n\ge1}(A[1])^{\otimes n}$ 上给定一个 $+1$ 次的
余微分（codifferential）$d$。
\end{dfn}

按定义，余微分 $d$ 是一个余导子，且满足 $d^2=0$。它由它的「Taylor 系数」
$m_n\colon A^{\otimes n}\to A[2-n]$（$n\ge1$）唯一确定。条件 $d^2=0$ 可以改写为一列二次方程
\[
\sum_{i+j=n+1}\ \sum_{0\le l\le i}\epsilon(l,j)\,
m_i(a_0,\dots,a_{l-1},m_j(a_l,\dots,a_{l+j-1}),a_{l+j},\dots,a_n)=0,
\]
其中 $a_m\in A$，而
\[
\epsilon(l,j)=(-1)^{\,j\sum_{0\le s\le l-1}\deg(a_s)+l(j-1)+j(i-1)}.
\]
特别地，$m_1^2=0$。

\begin{yizhu}[余代数表述与 Koszul 符号]
把一族运算 $\{m_n\}_{n\ge1}$ 打包成一个余导子 $d$，是 $A_\infty$-代数的标准表述：
$d$ 在张量幂 $(A[1])^{\otimes n}$ 上的分量就是 $m_n$（连同平移带来的符号）。方程
$d^2=0$ 展开后即得上面的二次方程组；$\epsilon(l,j)$ 就是 Koszul 符号，其来源是把
$m_j$ 移过 $a_0,\dots,a_{l-1}$ 时逐次交换所产生的符号，再加上 $m_j$ 与 $a_l$ 各自的
平移次数。核对方程中的符号时，按「先数被移过的元素的次数、再数平移次数」的顺序计算即可。
\end{yizhu}

\begin{dfn}
非幺 $A_\infty$-代数 $(V,d_V)$ 与 $(W,d_W)$ 之间的态射（$A_\infty$-态射）是指张量余代数
之间的一个零次态射 $T_+(V[1])\to T_+(W[1])$，它与各自的余微分交换。
\end{dfn}

非幺 $A_\infty$-代数的态射 $f$ 由它的「Taylor 系数」$f_n\colon V^{\otimes n}\to W[1-n]$
（$n\ge1$）确定，这些系数满足方程组
\[
\begin{aligned}
&\sum_{1\le l_1<\cdots<l_i=n}(-1)^{\gamma_i}m_i^W\bigl(f_{l_1}(a_1,\dots,a_{l_1}),\\
&\qquad\qquad f_{l_2-l_1}(a_{l_1+1},\dots,a_{l_2}),\dots,
 f_{n-l_{i-1}}(a_{n-l_{i-1}+1},\dots,a_n)\bigr)\\
&\quad=\sum_{s+r=n+1}\ \sum_{1\le j\le s}(-1)^{\epsilon_s}
f_s\bigl(a_1,\dots,a_{j-1},m_r^V(a_j,\dots,a_{j+r-1}),a_{j+r},\dots,a_n\bigr).
\end{aligned}
\]
这里
$\epsilon_s=r\sum_{1\le p\le j-1}\deg(a_p)+j-1+r(s-j)$，
\[
\gamma_i=\sum_{1\le p\le i-1}(i-p)(l_p-l_{p-1}-1)
+\sum_{1\le p\le i-1}\nu(l_p)\sum_{l_{p-1}+1\le q\le l_p}\deg(a_q),
\]
其中用到记号 $\nu(l_p)=\sum_{p+1\le m\le i}(1-l_m+l_{m-1})$，并约定 $l_0=0$。

\begin{dfn}
域 $k$ 上的非幺 $A_\infty$-范畴 $\mathcal{C}$ 由以下数据给出：
\begin{enumerate}[label=\arabic*),leftmargin=2.2em,itemsep=1pt,topsep=2pt]
\item 一个对象类 $\Ob(\mathcal{C})$。
\item 对任意两个对象 $X_1$ 与 $X_2$，一个 $\Z$-分次 $k$-模 $\Hom(X_1,X_2)$。
\item 对任意对象序列 $X_0,\dots,X_n$（$n\ge1$），一个 $k$-模态射（称为复合映射）
  \[
  \textstyle m_n\colon\bigotimes_{0\le i\le n-1}\Hom(X_i,X_{i+1})
  \longrightarrow\Hom(X_0,X_n)[2-n].
  \]
\end{enumerate}
要求：对任意对象序列 $X_0,\dots,X_N$（$N\ge0$），分次 $k$-模
$A=A(X_0,\dots,X_N):=\bigoplus_{i,j}\Hom(X_i,X_j)$ 配以复合 $m_n$（$n\ge1$）的直和，
构成一个非幺 $A_\infty$-代数。
\end{dfn}

对象类 $\Ob(\mathcal{C})$ 今后常简记为 $\mathcal{C}$。希望这不会引起混淆。

\begin{rmk}
非幺 $A_\infty$-代数 $A$ 可以看成一个非幺 $A_\infty$-范畴，它只有一个对象 $X$，且
$\Hom(X,X)=A$。
\end{rmk}

\begin{dfn}
非幺 $A_\infty$-范畴之间的函子 $F\colon\mathcal{C}_1\to\mathcal{C}_2$ 由以下数据给出：
\begin{enumerate}[label=\arabic*),leftmargin=2.2em,itemsep=1pt,topsep=2pt]
\item 对象类之间的映射 $\phi\colon\mathcal{C}_1\to\mathcal{C}_2$。
\item 对任意有限对象序列 $X_0,\dots,X_n$（$n\ge0$），一个分次 $k$-模的态射
  \[
  \textstyle f_n\colon\bigotimes_{0\le i\le n-1}\Hom_{\mathcal{C}_1}(X_i,X_{i+1})
  \longrightarrow\Hom_{\mathcal{C}_2}(\phi(X_0),\phi(X_n))[1-n].
  \]
\end{enumerate}
以下条件成立：对任意 $X_1,\dots,X_N\in\mathcal{C}_1$，序列 $f_n$（$n\ge1$）定义一个
$A_\infty$-态射
\[
\bigoplus_{i,j}\Hom_{\mathcal{C}_1}(X_i,X_j)\longrightarrow
\bigoplus_{i,j}\Hom_{\mathcal{C}_2}(\phi(X_i),\phi(X_j)).
\]
\end{dfn}

\begin{rmk}
设 $\mathcal{C}$ 是非幺 $A_\infty$-范畴。把态射空间替换为它关于 $m_1$ 的上同调。换言之，
定义 $\Hom_{H(\mathcal{C})}(X,Y):={\ker m_1}/{\im m_1}$，其中 $m_1$ 是复合映射
$\Hom_{\mathcal{C}}(X,Y)\to\Hom_{\mathcal{C}}(X,Y)[1]$。于是
$H(\mathcal{C})=(\mathcal{C},\Hom_{H(\mathcal{C})}(\cdot,\cdot))$ 给出一个「非幺」范畴
结构，其对象类为 $\mathcal{C}$，态射的复合由 $m_2$ 诱导。这里说「非幺」，是因为没有
单位态射 $\id_X\in\Hom_{H(\mathcal{C})}(X,X)$。
\end{rmk}

\subsection{$A_\infty$-预范畴}

我们先引进非幺 $A_\infty$-预范畴的概念。它使我们可以对「横截」的对象序列工作。\footnote{本节
的「横截性」概念纯属形式。这一名称的含义将在几何背景下的具体应用之后变得清楚，见后面几节。}
然后我们将引进 $A_\infty$-预范畴的概念。它为我们提供单位态射的替代物。粗略地说，我们将有
同伦意义下的单位态射。

\begin{dfn}
设 $k$ 是如前所述的 $\Z$-分次交换结合环。$k$ 上的非幺 $A_\infty$-预范畴由以下数据定义：
\begin{enumerate}[label=\alph*),leftmargin=2.2em,itemsep=1pt,topsep=2pt]
\item 一个对象类 $\mathcal{C}$。
\item 对任意 $n\ge1$，一个子类 $\mathcal{C}^{\mathrm{tr}}_n\subset\mathcal{C}^n$，
  其中 $\mathcal{C}^{\mathrm{tr}}_1=\mathcal{C}$，称为横截序列类。
\item 对 $(X_1,X_2)\in\mathcal{C}^{\mathrm{tr}}_2$，一个 $\Z$-分次 $k$-模
  $\Hom(X_1,X_2)$。
\item 对横截的对象序列 $(X_0,\dots,X_n)$（$n\ge0$），一个 $k$-模态射（复合映射）
  \[
  \textstyle m_n\colon\bigotimes_{0\le i\le n-1}\Hom(X_i,X_{i+1})
  \longrightarrow\Hom(X_0,X_n)[2-n].
  \]
\end{enumerate}
要求：横截序列 $(X_1,\dots,X_n)$ 的任意子序列 $(X_{i_1},\dots,X_{i_l})$
（$i_1<i_2<\cdots<i_l$）都是横截的，且复合映射满足与非幺 $A_\infty$-范畴相同的一组方程。
具体写出即
\[
\sum_{i+j=n+1}\ \sum_{0\le l\le i}\epsilon(l,j)\,
m_i(a_0,\dots,a_{l-1},m_j(a_l,\dots,a_{l+j-1}),a_{l+j},\dots,a_n)=0,
\]
其中 $a_m\in\Hom(X_m,X_{m+1})$，而
$\epsilon(l,j)=(-1)^{\,j\sum_{0\le s\le l-1}\deg(a_s)+l(j-1)+j(i-1)}$。
\end{dfn}

\begin{yizhu}[原文的一处下标笔误]
原文在此处把 $m_j$ 的末一个变量写作 $a_{l+j}$、把紧随其后的变量写作 $a_{l+j+1}$，与
$m_j$ 的元数（$j$ 个变量）以及定义 3 中的同一方程都不合。译文按数学意义改为 $a_{l+j-1}$
与 $a_{l+j}$。
\end{yizhu}

\begin{dfn}
非幺 $A_\infty$-预范畴之间的函子 $F\colon\mathcal{C}\to\mathcal{D}$ 由以下数据给出：
\begin{enumerate}[label=\arabic*),leftmargin=2.2em,itemsep=1pt,topsep=2pt]
\item 对象类的映射 $\phi\colon\mathcal{C}\to\mathcal{D}$，满足
  $\phi^n(\mathcal{C}^{\mathrm{tr}}_n)\subset\mathcal{D}^{\mathrm{tr}}_n$。
\item 对 $\mathcal{C}$ 中任意横截对象序列 $(X_0,\dots,X_n)$（$n\ge1$），一个分次 $k$-模
  态射
  \[
  \textstyle f_n\colon\bigotimes_{0\le i\le n-1}\Hom_{\mathcal{C}}(X_i,X_{i+1})
  \longrightarrow\Hom_{\mathcal{D}}(\phi(X_0),\phi(X_n))[1-n].
  \]
\end{enumerate}
这些数据满足以下性质：序列 $f_n$（$n\ge1$）定义一个 $A_\infty$-态射
$\bigoplus_{i<j}\Hom_{\mathcal{C}}(X_i,X_j)\to
\bigoplus_{i<j}\Hom_{\mathcal{D}}(\phi(X_i),\phi(X_j))$。
\end{dfn}

读者已经注意到，我们只对递增的指标对 $i<j$ 求和。这与非幺 $A_\infty$-范畴的情形不同。
原因是，我们不要求横截性是关于对象的对称关系。可能 $\Hom(X_0,X_1)$ 存在，而
$\Hom(X_1,X_0)$ 不存在。当所有 $\Hom$ 都有定义时，所讨论的两个定义一致。特别地，非幺
$A_\infty$-范畴等同于满足 $\mathcal{C}^{\mathrm{tr}}_n=\mathcal{C}^n$（对一切 $n\ge1$）
的非幺 $A_\infty$-预范畴。

\begin{yizhu}[原文的一处笔误]
原文此句作「This differs from the case of non-unital $A_\infty$-pre-categories」，
但比较的对象是定义 6 与定义 8 中函子所诱导的 $A_\infty$-态射：前者对 $\bigoplus_{i,j}$
求和，后者只对 $\bigoplus_{i<j}$ 求和。故此处应为「非幺 $A_\infty$-范畴」，译文已按此译出。
\end{yizhu}

\begin{dfn}
设 $\mathcal{C}$ 是非幺 $A_\infty$-预范畴，$(X_1,X_2)\in\mathcal{C}^{\mathrm{tr}}_2$。
称 $f\in\Hom^0(X_1,X_2)$（上标 $0$ 表示次数）为拟同构（quasi-isomorphism），若
$m_1(f)=0$，且对任意满足 $(X_0,X_1,X_2)\in\mathcal{C}^{\mathrm{tr}}_3$ 与
$(X_1,X_2,X_3)\in\mathcal{C}^{\mathrm{tr}}_3$ 的对象 $X_0$ 与 $X_3$，都有：
$m_2(f,\cdot)\colon\Hom(X_0,X_1)\to\Hom(X_0,X_2)$ 与
$m_2(\cdot,f)\colon\Hom(X_2,X_3)\to\Hom(X_1,X_3)$ 是复形的拟同构。
\end{dfn}

\begin{dfn}
$A_\infty$-预范畴是指满足以下扩张性质（extension property）的非幺 $A_\infty$-预范畴
$\mathcal{C}$：

对 $\mathcal{C}$ 中任意有限族横截序列 $(S_i)_{i\in I}$ 与任意对象 $X$，存在对象 $X_+$
与 $X_-$，以及拟同构 $f_-\colon X_-\to X$、$f_+\colon X\to X_+$，使得对一切 $i\in I$，
扩充后的序列 $(X_-,S_i,X_+)$ 都是横截的。
\end{dfn}

\begin{yizhu}[扩张性质的几何来源（译者补充）]
扩张性质把「单位态射」换成「在序列两端各接上一个与 $X$ 拟同构的对象」：$X_\pm$ 相当于
$X$ 的扰动，$f_\pm$ 记录扰动与原对象之间的拟同构。在 §5 的几何实现中，这正是用哈密顿同痕
（Hamiltonian isotopy）把拉格朗日子流形微扰成横截位置的标准手法，而扰动前后的对象由
Floer 同调中的拟同构相联系。原文此处未点明这一来源，本注为译者补充。
\end{yizhu}

\begin{rmk}
设 $\mathcal{C}$ 是 $A_\infty$-预范畴。则在
$H(\mathcal{C})=(\mathcal{C},\Hom_{H(\mathcal{C})}(\cdot,\cdot))$ 上部分定义的复合
$m_2$ 可以唯一地延拓，从而在 $H(\mathcal{C})$ 上定义一个范畴结构。
\end{rmk}

\begin{dfn}
设 $\mathcal{C}$ 与 $\mathcal{D}$ 是 $k$ 上的 $A_\infty$-预范畴。$A_\infty$-函子
$F\colon\mathcal{C}\to\mathcal{D}$ 是指相应的非幺 $A_\infty$-预范畴之间的函子，且 $F$
把 $\mathcal{C}$ 中的拟同构映为 $\mathcal{D}$ 中的拟同构。
\end{dfn}

$A_\infty$-预范畴（以及 $A_\infty$-范畴）的等价有一个重要的概念。我们计划在别处给出全部
细节（见 \cite{KoS}）。就本文的目的而言，我们将使用下面的定义（它在更一般的框架中其实
是一个定理）。

\begin{dfn}
称 $A_\infty$-预范畴之间的 $A_\infty$-函子 $F\colon\mathcal{C}\to\mathcal{D}$ 为
$A_\infty$-等价函子，若：
\begin{enumerate}[label=\alph*),leftmargin=2.2em,itemsep=1pt,topsep=2pt]
\item $\mathcal{D}$ 中每个对象 $Y$ 都拟同构于某个对象 $\phi(X)$，$X\in\mathcal{C}$。
\item 该函子对一切横截的对象序列，诱导出相应态射的非幺 $A_\infty$-代数之间的拟同构。
\end{enumerate}
\end{dfn}

\begin{dfn}
称两个 $A_\infty$-预范畴 $\mathcal{C}$ 与 $\mathcal{D}$ \emph{等价}，若存在有限的
$A_\infty$-预范畴序列 $(\mathcal{C}_0,\dots,\mathcal{C}_n)$，$\mathcal{C}_0=\mathcal{C}$、
$\mathcal{C}_n=\mathcal{D}$，使得对每个 $i$（$0\le i\le n-1$）都存在从 $\mathcal{C}_i$
到 $\mathcal{C}_{i+1}$ 或反方向的 $A_\infty$-等价函子。
\end{dfn}

\begin{yizhu}[原文的一处笔误]
原文此处作「$\mathcal{C}_0=\mathcal{C}$，$\mathcal{C}_0=\mathcal{D}$」以及
「$0\le i\le k-1$」，与所取的有限序列 $(\mathcal{C}_0,\dots,\mathcal{C}_n)$ 不合。
按上下文，应为 $\mathcal{C}_0=\mathcal{C}$、$\mathcal{C}_n=\mathcal{D}$ 与
$0\le i\le n-1$。译文已按此译出。
\end{yizhu}

我们建议使用 $A_\infty$-预范畴的语言，以替代带严格单位态射（strict identity morphism）的、
更传统的 $A_\infty$-范畴。

\begin{dfn}
带严格单位态射的 $A_\infty$-范畴是指这样的非幺 $A_\infty$-范畴 $\mathcal{C}$：对任意对象
$X$，存在元素 $1=1_X\in\Hom^0(X,X)$（单位态射），使得 $m_2(1,f)=m_2(f,1)=f$，且对任意
态射 $f,f_1,\dots,f_n$ 有 $m_n(f_1,\dots,1,\dots,f_n)=0$（$n\ne2$）。
\end{dfn}

带严格单位态射的 $A_\infty$-范畴是 $A_\infty$-预范畴，因为（沿用前面的记号）我们可以把
横截序列 $S$ 扩充为 $(X,S,X)$，并取 $X_+=X_-=X$、$f_\pm=1_X$。另一点要说明的是，若
$\mathcal{C}$ 只有一个对象，它就是带严格单位的 $A_\infty$-代数。可以沿 \cite{KoS1} 的
思路发展这类代数的形变理论。问题在于相应的 operad 不是自由的，标准理论因此变得复杂。
我们希望 $A_\infty$-预范畴的框架适用于 $A_\infty$-范畴的形变理论。下面的猜想为
$A_\infty$-范畴的这种推广提供了另一佐证。

\begin{conj}
我们类似 $A_\infty$-预范畴的情形（见上）来定义带严格单位态射的 $A_\infty$-范畴的等价。
于是 $A_\infty$-预范畴的等价类与带严格单位态射的 $A_\infty$-范畴的等价类一一对应。
\end{conj}

\subsection{例：有向 $A_\infty$-预范畴}

$A_\infty$-预范畴的概念有一个有用的特殊情形（一个类似的概念由 \cite{Se} 独立提出）。

\begin{dfn}
有向 $A_\infty$-预范畴（directed $A_\infty$-pre-category）是指满足以下条件的
$A_\infty$-预范畴：
\begin{enumerate}[label=\alph*),leftmargin=2.2em,itemsep=1pt,topsep=2pt]
\item 对象类与整数集之间存在双射：$\mathcal{C}\simeq\Z$。用 $X_i$ 记对应于 $i\in\Z$
  的对象。
\item 横截序列是 $(X_{i_1},\dots,X_{i_n})$，其中 $i_1<i_2<\cdots<i_n$。
\end{enumerate}
\end{dfn}

扩张性质等价于下述性质：对任意对象 $X_i$，存在对象 $X_j$（$j<i$）与 $X_m$（$m>i$），
它们都拟同构于 $X_i$。于是可以陈述前述猜想的如下版本。

\begin{conj}
有向 $A_\infty$-预范畴的等价类，与带严格单位态射且对象类可数的 $A_\infty$-范畴的等价类
一一对应。
\end{conj}

给定一个带严格单位态射、对象类可数的 $A_\infty$-范畴 $\mathcal{C}$，可以构造一个无穷对象
序列 $(X_i)_{i\in\Z}$，使每个对象在正指标与负指标处都出现无穷多次。然后如下定义有向
$A_\infty$-预范畴 $\mathcal{C}'$：令 $\Hom_{\mathcal{C}'}(X_i,X_j)=
\Hom_{\mathcal{C}}(X_i,X_j)$（$i<j$），其余 $\Hom$ 均不定义。

% [原文 p22–p28；§4.4 止于 p28，§5 自 p28 起]
